\documentclass[11pt,a4paper]{article}
\usepackage[utf8]{inputenc}
\usepackage[italian]{babel}
\usepackage{geometry}
\usepackage{graphicx}
\usepackage{hyperref}
\usepackage{listings}
\usepackage{xcolor}
\usepackage{amsmath}
\usepackage{enumitem}
\usepackage{booktabs}
\usepackage{longtable}
\usepackage{array}
\usepackage{tcolorbox}
\usepackage{float}

\geometry{margin=2.5cm}

% Configurazione codice
\lstset{
    language=JavaScript,
    basicstyle=\ttfamily\small,
    keywordstyle=\color{blue}\bfseries,
    stringstyle=\color{red},
    commentstyle=\color{green!60!black},
    numberstyle=\tiny\color{gray},
    numbers=left,
    numbersep=5pt,
    frame=single,
    breaklines=true,
    showstringspaces=false,
    tabsize=2,
    escapeinside={(*@}{@*)}
}

% Box per evidenziare sezioni importanti
\newtcolorbox{importantbox}[1]{
    colback=red!5!white,
    colframe=red!75!black,
    title=#1,
    fonttitle=\bfseries
}

\newtcolorbox{warningbox}[1]{
    colback=orange!5!white,
    colframe=orange!75!black,
    title=#1,
    fonttitle=\bfseries
}

\newtcolorbox{tipbox}[1]{
    colback=blue!5!white,
    colframe=blue!75!black,
    title=#1,
    fonttitle=\bfseries
}

\newtcolorbox{successbox}[1]{
    colback=green!5!white,
    colframe=green!75!black,
    title=#1,
    fonttitle=\bfseries
}

\title{\textbf{WorkLogService: Documentazione Post-Refactoring}\\
\large Sistema Unificato con Event-Driven Architecture}
\author{Documentazione Tecnica - Versione Definitiva}
\date{\today}

\begin{document}

\maketitle
\tableofcontents
\newpage

\section{Introduzione}

\begin{importantbox}{Stato Attuale - Post-Refactoring Completo}
WorkLogService è il \textbf{Cuore Operativo} del sistema di tracciamento lavoro dopo il refactoring di unificazione completo.
Questo componente gestisce sia log di lavoro con orari (\textbf{Timer}) che note di lavoro senza orari (\textbf{Journal}) in un'unica entità unificata.
Il sistema utilizza un'architettura \textbf{Event-Driven} con Pattern Observer per decoupling e resilienza.
\end{importantbox}

\subsection{Contesto Storico e Unificazione}

Prima del refactoring, il sistema aveva entità separate che causavano duplicazione e complessità:
\begin{itemize}
    \item \textbf{WorkLog}: Solo per tracciamento orari (startTime, endTime, durationMinutes)
    \item \textbf{WorkEntry}: Solo per note di lavoro (title, description, tags, mood)
    \item \textbf{WorkProject}: Entità separata per progetti
\end{itemize}

Dopo il refactoring completo:
\begin{itemize}
    \item \textbf{WorkLog}: Entità \textbf{ibrida} che supporta sia Timer che Journal
    \item \textbf{Project}: Unica entità per progetti (eliminato WorkProject)
    \item \textbf{Pattern Observer}: Decoupling completo tramite Event Bus
    \item \textbf{Bull Queue}: Retry automatico per attività secondarie
\end{itemize}

\subsection{Architettura Unificata: Single Source of Truth}

Il sistema ora segue il principio \textbf{Single Source of Truth}:

\begin{itemize}
    \item \textbf{WorkLog} è l'unica entità per tracciamento lavoro
    \item \textbf{Project} è l'unica entità per progetti (con tipo CLIENT/PERSONAL)
    \item \textbf{Event Bus} centralizza la comunicazione tra servizi
    \item \textbf{Activity Queue} gestisce retry e resilienza
\end{itemize}

\section{Architettura del Sistema}

\subsection{Modello WorkLog: Entità Ibrida}

Il modello \texttt{WorkLog} supporta due modalità di utilizzo:

\subsubsection{Modalità Timer (Time Tracking)}
Quando \texttt{startTime} e \texttt{endTime} sono presenti:
\begin{itemize}
    \item La durata viene calcolata automaticamente (salvo lock manuale)
    \item Supporta lavoro oltre mezzanotte (es. 23:00 - 02:00)
    \item Genera attività di tipo \texttt{WORK\_SESSION\_LOGGED}
\end{itemize}

\subsubsection{Modalità Journal (Note)}
Quando \texttt{startTime} e \texttt{endTime} sono assenti:
\begin{itemize}
    \item \texttt{durationMinutes} rimane 0
    \item Permette note di lavoro senza tracciamento orario
    \item Genera attività di tipo \texttt{WORK\_NOTE\_CREATED}
\end{itemize}

\begin{lstlisting}[caption=Schema WorkLog Unificato]
const workLogSchema = new mongoose.Schema({
    projectId: { type: ObjectId, ref: 'Project', required: true },
    date: { type: String, required: true, match: /^\d{4}-\d{2}-\d{2}$/ },
    title: { type: String, required: true, maxlength: 200 },
    description: { type: String, maxlength: 5000 },
    tags: [{ type: String, maxlength: 50 }],
    mood: { type: String, enum: ['high', 'neutral', 'low'] },
    isBillable: { type: Boolean, default: true },
    
    // Timer fields (opzionali)
    startTime: { type: String, match: /^\d{2}:\d{2}$/ },
    endTime: { type: String, match: /^\d{2}:\d{2}$/ },
    durationMinutes: { type: Number, min: 0, default: 0 },
    
    // Lock durata manuale
    isManualDuration: { type: Boolean, default: false },
}, { timestamps: true });
\end{lstlisting}

\subsection{Modello Project: Tipi e Validazione}

Il modello \texttt{Project} supporta due tipi con validazione differenziata:

\begin{lstlisting}[caption=Schema Project con Tipi]
const projectSchema = new mongoose.Schema({
    name: { type: String, required: true, maxlength: 100 },
    code: { type: String, required: true, uppercase: true },
    description: { type: String, maxlength: 500 },
    
    // Tipo progetto: CLIENT o PERSONAL
    type: {
        type: String,
        enum: ['CLIENT', 'PERSONAL'],
        default: 'CLIENT',
    },
    
    // Rate: obbligatorio per CLIENT, opzionale per PERSONAL
    rate: {
        type: Number,
        required: function() {
            return this.type === 'CLIENT';
        },
        min: 0,
    },
    
    estimatedHours: { type: Number, min: 0, default: 0 },
    progress: { type: Number, min: 0, max: 100, default: 0 },
    status: { type: String, enum: ['active', 'completed', 'archived'] },
    budget: { type: Number, min: 0 }, // Solo per CLIENT
    color: { type: String, default: '#6366f1' },
}, { timestamps: true });
\end{lstlisting}

\begin{tipbox}{Validazione Rate per Tipo}
\begin{itemize}
    \item \textbf{CLIENT}: \texttt{rate} è \textbf{obbligatorio} (pre-save hook valida)
    \item \textbf{PERSONAL}: \texttt{rate} è \textbf{opzionale} (default: 0)
    \item Il frontend nasconde/mostra campi budget in base al tipo
\end{itemize}
\end{tipbox}

\section{Logica di Business Avanzata}

\subsection{Durata Manuale: Lock e Pre-Save Hook}

Il sistema supporta la modifica manuale della durata per gestire pause non tracciate.

\begin{importantbox}{Campo isManualDuration}
Quando l'utente modifica manualmente \texttt{durationMinutes} nel frontend, viene impostato \texttt{isManualDuration: true}.
Il pre-save hook rispetta questo flag e \textbf{non ricalcola} la durata anche se \texttt{startTime}/\texttt{endTime} sono presenti.
\end{importantbox}

\begin{lstlisting}[caption=Pre-Save Hook con Lock Durata]
workLogSchema.pre('save', function(next) {
    // Se la durata è manuale, salta il calcolo automatico
    if (this.isManualDuration) {
        return next();
    }

    // Calcola durationMinutes solo se startTime E endTime sono presenti
    if (this.startTime && this.endTime) {
        if (this.isModified('startTime') || this.isModified('endTime')) {
            const [startHour, startMin] = this.startTime.split(':').map(Number);
            const [endHour, endMin] = this.endTime.split(':').map(Number);
            
            let duration = (endHour * 60 + endMin) - (startHour * 60 + startMin);
            
            // Gestisce lavoro oltre mezzanotte
            if (duration < 0) {
                duration += 24 * 60;
            }
            
            this.durationMinutes = duration;
        }
    } else {
        // Se orari mancano, durationMinutes = 0
        if (!this.isModified('durationMinutes')) {
            this.durationMinutes = 0;
        }
    }
    
    next();
});
\end{lstlisting}

\subsubsection{UX Frontend: Lock Manuale}

Nel frontend (\texttt{WorkLogFormSmart.tsx}), quando l'utente modifica manualmente la durata:
\begin{itemize}
    \item Si attiva automaticamente il lock (\texttt{isManualDuration: true})
    \item Viene mostrato un indicatore visivo (icona lock, bordo colorato)
    \item L'utente può sbloccare per tornare al calcolo automatico
    \item Modifiche a \texttt{startTime}/\texttt{endTime} non ricalcolano se lock attivo
\end{itemize}

\subsection{Gestione Mezzanotte}

Il calcolo automatico della durata gestisce correttamente il passaggio di giorno:

\begin{lstlisting}[caption=Calcolo Durata con Mezzanotte]
let duration = (endHour * 60 + endMin) - (startHour * 60 + startMin);

// Gestisce lavoro oltre mezzanotte
// Esempio: 23:00 - 02:00 = 3 ore (non -21 ore)
if (duration < 0) {
    duration += 24 * 60; // Aggiungi 24 ore
}
\end{lstlisting}

\begin{tipbox}{Esempi Mezzanotte}
\begin{itemize}
    \item \texttt{23:00 - 02:00} = 3 ore (180 minuti)
    \item \texttt{22:00 - 06:30} = 8.5 ore (510 minuti)
    \item \texttt{18:00 - 02:30} = 8.5 ore (510 minuti)
\end{itemize}
\end{tipbox}

\subsection{Validazione Ownership Progetto}

Ogni operazione su WorkLog verifica che il progetto appartenga allo stesso utente:

\begin{lstlisting}[caption=Validazione Ownership Atomica]
async validateProjectOwnership(tenantScope, projectId) {
    const userId = this._getUserId(tenantScope);
    
    // Query diretta con filtro esplicito (IDOR prevention)
    const project = await Project.findOne({ 
        _id: projectId, 
        user: userId 
    });
    
    if (!project) {
        // Non rivelare se progetto esiste ma appartiene ad altri
        throw AppError.notFound('Progetto');
    }
}
\end{lstlisting}

\begin{warningbox}{Sicurezza: IDOR Prevention}
Senza questa validazione, un utente malintenzionato potrebbe:
\begin{enumerate}
    \item Conoscere l'ID di un progetto altrui (da URL o leak)
    \item Creare un WorkLog associato a quel progetto
    \item "Inquinare" i dati di un altro utente
\end{enumerate}
La validazione è implementata in \texttt{beforeCreate()} e \texttt{beforeUpdate()}.
\end{warningbox}

\section{Event-Driven Architecture}

\subsection{Pattern Observer: Decoupling Completo}

Il sistema utilizza il \textbf{Pattern Observer} per disaccoppiare i servizi di dominio dall'activity tracking.

\begin{figure}[H]
\centering
\begin{lstlisting}[caption=Architettura Event-Driven]
┌─────────────────┐
│  WorkLogService │
│  GoalService    │  ──emit──>  ┌──────────┐
│  StudyService   │             │ EventBus │
└─────────────────┘             └────┬─────┘
                                     │
                                     │ on()
                                     ▼
                            ┌─────────────────┐
                            │ ActivitySubscriber│
                            └────┬────────────┘
                                 │
                                 │ addActivityJob()
                                 ▼
                            ┌─────────────────┐
                            │  ActivityQueue  │
                            │  (Bull + Redis) │
                            └────┬────────────┘
                                 │
                                 │ retry (3x)
                                 ▼
                            ┌─────────────────┐
                            │ ActivityService │
                            └─────────────────┘
\end{lstlisting}
\caption{Flusso Event-Driven Architecture}
\end{figure}

\subsection{Event Bus: Singleton Centralizzato}

L'Event Bus è un singleton che gestisce tutti gli eventi dell'applicazione:

\begin{lstlisting}[caption=Event Bus Implementation]
class AppEventBus extends EventEmitter {
    constructor() {
        super();
        this.setMaxListeners(20);
    }

    emit(event, data) {
        // Emissione asincrona (non blocca call stack)
        setImmediate(() => {
            try {
                super.emit(event, data);
            } catch (err) {
                console.error(`❌ EventBus error (${event}):`, err.message);
            }
        });
    }
}

const eventBus = new AppEventBus();
\end{lstlisting}

\subsubsection{Eventi Emessi da WorkLogService}

\begin{itemize}
    \item \texttt{worklog.created}: Quando viene creato un nuovo worklog
    \begin{lstlisting}[caption=Payload worklog.created]
{
    userId: ObjectId,
    workLog: {
        _id: ObjectId,
        projectId: ObjectId,
        date: string,
        startTime: string | null,
        endTime: string | null,
        durationMinutes: number,
        tags: string[],
        // ... altri campi
    }
}
    \end{lstlisting}
    
    \item \texttt{worklog.updated}: Quando viene aggiornato un worklog
    \begin{lstlisting}[caption=Payload worklog.updated]
{
    userId: ObjectId,
    workLog: WorkLog,        // Documento aggiornato
    previousWorkLog: WorkLog // Documento precedente
}
    \end{lstlisting}
\end{itemize}

\subsection{Activity Queue: Retry con Exponential Backoff}

La queue Bull gestisce il retry automatico per attività secondarie:

\begin{lstlisting}[caption=Configurazione Activity Queue]
activityQueue = new Queue('activity-tracking', {
    redis: process.env.REDIS_URL,
    defaultJobOptions: {
        attempts: 3,              // 3 tentativi totali
        backoff: {
            type: 'exponential',  // Exponential backoff
            delay: 2000,          // Delay iniziale: 2s
        },
        removeOnComplete: {
            age: 3600,            // Rimuovi dopo 1 ora
            count: 1000,          // Max 1000 job completati
        },
    },
});
\end{lstlisting}

\subsubsection{Retry Strategy}

\begin{itemize}
    \item \textbf{Tentativo 1}: Immediato
    \item \textbf{Tentativo 2}: Dopo 2 secondi (exponential backoff)
    \item \textbf{Tentativo 3}: Dopo 4 secondi (exponential backoff)
    \item \textbf{Fallimento}: Dopo 3 tentativi, job marcato come fallito
\end{itemize}

\subsubsection{Fallback}

Se Redis non è disponibile:
\begin{itemize}
    \item La queue non viene inizializzata
    \item \texttt{addActivityJob()} usa chiamata diretta a \texttt{ActivityService}
    \item Nessun retry automatico (ma non blocca l'app)
\end{itemize}

\subsection{Activity Subscriber: Handler Eventi}

Il subscriber ascolta gli eventi e li processa tramite la queue:

\begin{lstlisting}[caption=Handler worklog.created]
async function handleWorkLogCreated(data) {
    const { userId, workLog } = data;
    
    const activityType = (workLog.startTime && workLog.endTime) 
        ? 'WORK_SESSION_LOGGED' 
        : 'WORK_NOTE_CREATED';

    // Usa queue con retry automatico
    await addActivityJob(userId, activityType, {
        entityId: workLog._id || workLog.id,
        category: 'work',
        metadata: {
            projectId: workLog.projectId,
            durationMinutes: workLog.durationMinutes || 0,
            tags: workLog.tags || [],
            timeOfDay: getTimeOfDayLabel(new Date()),
            date: workLog.date,
            startTime: workLog.startTime || null,
            endTime: workLog.endTime || null,
            hasTimeTracking: !!(workLog.startTime && workLog.endTime),
        },
    });
}
\end{lstlisting}

\subsection{Event Metrics: Monitoring}

Il sistema traccia metriche automaticamente per monitoring:

\begin{lstlisting}[caption=Event Metrics]
class EventMetrics {
    constructor() {
        this.metrics = {
            emitted: {},      // Eventi emessi per tipo
            processed: {},    // Eventi processati con successo
            failed: {},       // Eventi falliti
            processingTime: {}, // Tempo medio di processing
        };
    }
    
    getStats() {
        return {
            emitted: { ...this.metrics.emitted },
            processed: { ...this.metrics.processed },
            failed: { ...this.metrics.failed },
            successRate: {},        // Tasso di successo per tipo
            avgProcessingTime: {}, // Tempo medio per tipo
        };
    }
}
\end{lstlisting}

\begin{tipbox}{Metriche Disponibili}
\begin{itemize}
    \item Eventi emessi per tipo (worklog.created, worklog.updated, etc.)
    \item Eventi processati con successo
    \item Eventi falliti (con retry)
    \item Success rate per tipo di evento
    \item Tempo medio di processing
    \item Logging automatico ogni 5 minuti in sviluppo
\end{itemize}
\end{tipbox}

\section{Flusso Operativo Completo}

\subsection{Creazione WorkLog: Da UI a Database}

\begin{figure}[H]
\centering
\begin{lstlisting}[caption=Flusso Creazione WorkLog]
[Frontend] WorkLogFormSmart.tsx
     │
     │ onSubmit() → onSave({ 
     │   projectId, date, startTime, endTime,
     │   durationMinutes?, isManualDuration?
     │ })
     │
     ▼
[API] POST /api/worklogs
     │
     ▼
[Controller] workLogService.create(req.tenantScope, data)
     │
     ├─► [beforeCreate()]
     │   │
     │   ├─► validateProjectOwnership()
     │   │   └─► Verifica che projectId appartenga all'utente
     │   │
     │   └─► validateTimeRange() (se orari presenti)
     │       └─► Verifica coerenza startTime/endTime
     │
     ├─► [super.create()]
     │   └─► BaseService.create() → WorkLog.create()
     │       └─► [Model pre-save hook]
     │           ├─► Se isManualDuration: skip calcolo
     │           └─► Altrimenti: calcola durationMinutes
     │
     ├─► [eventBus.emit('worklog.created')]
     │   └─► Fire and forget (non blocca risposta)
     │
     └─► [return created] → HTTP 201
     
[Background] ActivitySubscriber
     │
     ├─► handleWorkLogCreated()
     │   └─► addActivityJob() → ActivityQueue
     │       └─► [Bull Queue] retry automatico
     │           └─► activityService.recordActivity()
\end{lstlisting}
\caption{Flusso Completo Creazione WorkLog}
\end{figure}

\subsection{Aggiornamento WorkLog}

\begin{lstlisting}[caption=Flusso Update WorkLog]
async update(tenantScope, id, data, options = {}) {
    // 1. Recupera worklog precedente
    const previous = await this.findById(tenantScope, id);
    
    // 2. Esegui update normale
    const updated = await super.update(tenantScope, id, data, options);
    
    // 3. Emetti evento (Pattern Observer)
    eventBus.emit('worklog.updated', {
        userId,
        workLog: updated,
        previousWorkLog: previous,
    });
    
    return updated;
}
\end{lstlisting}

Il subscriber \texttt{handleWorkLogUpdated()} registra attività solo se ci sono cambiamenti significativi (durata o progetto).

\section{Integrazione Frontend}

\subsection{Navigazione: Projects → Workspace}

Il flusso di navigazione dal frontend:

\begin{enumerate}
    \item \textbf{Pagina Progetti}: L'utente vede la lista progetti
    \item \textbf{Click "Gestisci"}: Su una ProjectCard
    \item \textbf{Navigazione}: \texttt{navigate(`/workspace?projectId=\${projectId}`)}
    \item \textbf{WorkspacePage}: Legge \texttt{projectId} da URL query params
    \item \textbf{Filtro Automatico}: Mostra solo worklog del progetto selezionato
\end{enumerate}

\begin{lstlisting}[caption=Navigazione Frontend]
// ProjectCard.tsx
const handleManageClick = (e: React.MouseEvent) => {
    e.stopPropagation();
    const projectId = project.id || (project as any)._id;
    if (projectId) {
        navigate(`/workspace?projectId=${projectId}`);
    }
};

// WorkspacePage.tsx
const [searchParams] = useSearchParams();
const urlProjectId = searchParams.get('projectId');

useEffect(() => {
    if (urlProjectId) {
        // Sincronizza selectedProject con URL
        setSelectedProject(urlProjectId);
    }
}, [urlProjectId]);
\end{lstlisting}

\subsection{Adattamento UI in Base al Tipo Progetto}

Il frontend si adatta dinamicamente in base al tipo di progetto:

\begin{itemize}
    \item \textbf{CLIENT}: Mostra campi budget, rate, revenue
    \item \textbf{PERSONAL}: Nasconde campi budget, rate opzionale
    \item \textbf{Validazione}: Rate obbligatorio solo per CLIENT
\end{itemize}

\section{API Endpoints}

\subsection{GET /api/worklogs}

Lista tutti i worklog dell'utente con filtri opzionali.

\begin{lstlisting}[caption=Query Parameters]
GET /api/worklogs?projectId=123&date=2024-01-15&tags=urgent
\end{lstlisting}

\subsection{GET /api/worklogs/feed}

Feed Workspace con filtri avanzati (usato dalla vista operativa).

\begin{lstlisting}[caption=Feed Parameters]
GET /api/worklogs/feed?projectId=123&startDate=2024-01-01&endDate=2024-01-31&tags=urgent
\end{lstlisting}

\begin{itemize}
    \item \texttt{projectId}: Filtra per progetto specifico
    \item \texttt{date}: Filtra per data singola (YYYY-MM-DD)
    \item \texttt{startDate}: Inizio range date
    \item \texttt{endDate}: Fine range date
    \item \texttt{tags}: Filtra per tag (singolo o multipli)
    \item \texttt{limit}: Limite risultati (default: tutti)
    \item \texttt{skip}: Skip risultati (paginazione)
\end{itemize}

Ritorna worklog ordinati per \texttt{date} decrescente, poi \texttt{createdAt} decrescente (timeline-friendly).

\subsection{POST /api/worklogs}

Crea nuovo worklog (Timer o Journal).

\begin{lstlisting}[caption=Request Body]
{
    "projectId": "507f1f77bcf86cd799439011",
    "date": "2024-01-15",
    "title": "Sviluppo feature X",
    "description": "Implementato sistema di notifiche",
    "tags": ["frontend", "urgent"],
    "mood": "high",
    "isBillable": true,
    
    // Timer fields (opzionali)
    "startTime": "09:00",
    "endTime": "17:00",
    
    // Durata manuale (opzionale)
    "durationMinutes": 480,
    "isManualDuration": false
}
\end{lstlisting}

\subsection{PUT /api/worklogs/:id}

Aggiorna worklog esistente.

\subsection{GET /api/worklogs/by-month/:year/:month}

Worklog filtrati per mese specifico.

\begin{lstlisting}[caption=Esempio]
GET /api/worklogs/by-month/2024/01
\end{lstlisting}

\subsection{GET /api/worklogs/totals}

Totali ore e guadagni per un periodo.

\begin{lstlisting}[caption=Query Parameters]
GET /api/worklogs/totals?startDate=2024-01-01&endDate=2024-01-31
\end{lstlisting}

Ritorna:
\begin{itemize}
    \item \texttt{totalMinutes}: Minuti totali lavorati
    \item \texttt{totalEarnings}: Guadagni totali (ore × rate)
    \item \texttt{logCount}: Numero totale di log
\end{itemize}

\subsection{GET /api/worklogs/project/:projectId/stats}

Statistiche progetto per Vista Amministrativa (Dashboard Progetti).

\begin{lstlisting}[caption=Response]
{
    "projectId": "507f1f77bcf86cd799439011",
    "projectName": "Progetto X",
    "projectCode": "PROJ01",
    "totalHours": 120.5,
    "billableHours": 100.0,
    "nonBillableHours": 20.5,
    "totalRevenue": 5000.00,
    "burnRate": 75.5,  // Percentuale (null se estimatedHours = 0)
    "lastActivity": "2024-01-15",
    "entriesCount": 45,
    "projectRate": 50.0,
    "estimatedHours": 160.0
}
\end{lstlisting}

\begin{tipbox}{Burn Rate}
Il burn rate è calcolato come: \texttt{(ore consumate / ore stimate) × 100}.
Se \texttt{estimatedHours} è 0 o null, \texttt{burnRate} è \texttt{null} (non calcolabile).
\end{tipbox}

\section{Sicurezza e Multi-Tenancy}

\subsection{Isolamento Tenant}

Tutte le query includono automaticamente il filtro \texttt{user}:

\begin{lstlisting}[caption=Query Multi-Tenant]
// BaseService._buildFilter() aggiunge automaticamente:
const secureFilter = {
    _id: id,
    user: userId  // Sempre presente
};

const doc = await WorkLog.findOne(secureFilter);
\end{lstlisting}

\subsection{Validazione Ownership}

Ogni operazione che modifica un WorkLog verifica:
\begin{enumerate}
    \item Il worklog appartiene all'utente corrente
    \item Il progetto associato appartiene all'utente corrente
    \item Nessuna possibilità di IDOR (Insecure Direct Object Reference)
\end{enumerate}

\section{Performance e Ottimizzazioni}

\subsection{Indici Database}

\begin{lstlisting}[caption=Indici WorkLog]
// Query frequenti: worklog per utente e data
workLogSchema.index({ user: 1, date: -1 });

// Query frequenti: worklog per progetto
workLogSchema.index({ user: 1, projectId: 1, date: -1 });

// Ricerca testuale
workLogSchema.index({ title: 'text', description: 'text' });

// Ricerca per tag
workLogSchema.index({ tags: 1 });
\end{lstlisting}

\subsection{Fire and Forget}

Gli eventi vengono emessi in modo asincrono:
\begin{itemize}
    \item Non bloccano la risposta HTTP
    \item Processing in background tramite queue
    \item Retry automatico in caso di errori
\end{itemize}

\section{Testing e Debugging}

\subsection{Event Metrics in Sviluppo}

In ambiente di sviluppo, le metriche vengono loggate automaticamente ogni 5 minuti:

\begin{lstlisting}[caption=Log Metriche]
📊 Event Metrics: {
    total: {
        emitted: 150,
        processed: 148,
        failed: 2,
        successRate: "98.67%"
    },
    byEvent: {
        "worklog.created": {
            emitted: 100,
            processed: 99,
            failed: 1,
            successRate: "99.00%"
        }
    }
}
\end{lstlisting}

\subsection{Queue Stats}

Statistiche della queue Bull disponibili via API:

\begin{lstlisting}[caption=Queue Statistics]
const stats = await getQueueStats();
// {
//     waiting: 5,
//     active: 2,
//     completed: 100,
//     failed: 3,
//     delayed: 0,
//     total: 110
// }
\end{lstlisting}

\section{Conclusioni}

\subsection{Vantaggi Architettura Post-Refactoring}

\begin{enumerate}
    \item \textbf{Single Source of Truth}: WorkLog unificato elimina duplicazione
    \item \textbf{Decoupling}: Pattern Observer elimina accoppiamento stretto
    \item \textbf{Resilienza}: Retry automatico gestisce errori temporanei
    \item \textbf{Performance}: Fire and forget non blocca risposta HTTP
    \item \textbf{Scalabilità}: Facile aggiungere nuovi subscriber
    \item \textbf{Monitoring}: Metriche automatiche per debugging
\end{enumerate}

\subsection{Prossimi Sviluppi}

\begin{itemize}
    \item Webhook esterni per integrazioni
    \item Notifiche push per eventi importanti
    \item Dashboard analytics avanzata
    \item Export dati in formati multipli
\end{itemize}

\end{document}
